\documentclass{article}
\usepackage{graphicx}
\usepackage{pgfplots}
\usepackage{float}
\usepgfplotslibrary{fillbetween}
\pgfplotsset{compat=1.9}
\author{Emre Geçit}

\title{\includegraphics{9.4.png}\\ IE407 - Homework 1}

\begin{document}
\maketitle

\newpage
\section*{Question 1}
\paragraph*{a)}
As can be seen in Figure 1, optimal solution to this problem is on the point (200, 50), where the line for wood constraint and the cotton constraint intersect.
At this point, profit is 3200\$.
\newcommand\xmax{400}
\begin{figure}[H]
    \centering
    \begin{tikzpicture}
        \begin{axis}[
                legend style={nodes={scale=0.8}},
                grid,
                grid style={densely dashed},
                axis line style={->},
                xlabel=$x$,
                ylabel=$y$,
                xmin=0, xmax=\xmax,
                ymin=0, ymax=400,
                axis lines=center,
                axis on top=true,
                domain=0:\xmax,
            ]
            \addplot [draw=red,thick,name path=wool_constraint] {150-x/2}; \addlegendentry{Wool constraint: $3x + 6y \leq 900$}
            \addplot [draw=blue,thick,name path=cotton_constraint] {300-5*x/4}; \addlegendentry{Cotton constraint: $5x + 4y \leq 1200$}
            \addplot [draw=green, thick, name path=profit] {(3200-12*x)/16}; \addlegendentry{Profit: $12x + 16y = 3200$}
            \path[name path=xaxis] (0,0) -- (\xmax,0);
            \addplot[cyan!20] fill between[of=xaxis and wool_constraint,soft clip={domain=0:200}];
            \addplot[cyan!20] fill between[of=xaxis and cotton_constraint,soft clip={domain=200:\xmax}];
        \end{axis}
    \end{tikzpicture}
    \caption{Graphical Solution for Question 1a}
\end{figure}

\paragraph*{b)}
If the profit for a sweatshirt is increased by 1 and changed to 17\$, the profit function becomes $12x + 17y$.
As can be seen in Figure 2, the optimal solution for this problem does not change.
\begin{figure}[H]
    \centering
    \begin{tikzpicture}
        \begin{axis}[
                legend style={nodes={scale=0.8}},
                grid,
                grid style={densely dashed},
                axis line style={->},
                xlabel=$x$,
                ylabel=$y$,
                xmin=0, xmax=\xmax,
                ymin=0, ymax=400,
                axis lines=center,
                axis on top=true,
                domain=0:\xmax,
            ]
            \addplot [draw=red,thick,name path=wool_constraint] {150-x/2}; \addlegendentry{Wool constraint: $3x + 6y \leq 900$}
            \addplot [draw=blue,thick,name path=cotton_constraint] {300-5*x/4}; \addlegendentry{Cotton constraint: $5x + 4y \leq 1200$}
            \addplot [draw=green, thick, name path=profit] {(3200-12*x)/17}; \addlegendentry{Profit: $12x + 17y = 3200$}
            \path[name path=xaxis] (0,0) -- (\xmax,0);
            \addplot[cyan!20] fill between[of=xaxis and wool_constraint,soft clip={domain=0:200}];
            \addplot[cyan!20] fill between[of=xaxis and cotton_constraint,soft clip={domain=200:\xmax}];
        \end{axis}
    \end{tikzpicture}
    \caption{Graphical Solution for Question 1b}
\end{figure}

\paragraph*{c)}
When the profit for a t-shirt is changed to 20\$, the profit function becomes $20x + 16y$. In this case, the profit line's slope is same with the cotton constraint line.
The optimal solution is all the points on the cotton constraint line, which satisfy the wool constraint.

If the profit is increased even more, the wool constraint stops being a binding constraint, and the optimal solution changes to (240, 0).

\begin{figure}[H]
    \centering
    \begin{tikzpicture}
        \begin{axis}[
                legend style={nodes={scale=0.8}},
                grid,
                grid style={densely dashed},
                axis line style={->},
                xlabel=$x$,
                ylabel=$y$,
                xmin=0, xmax=\xmax,
                ymin=0, ymax=400,
                axis lines=center,
                axis on top=true,
                domain=0:\xmax,
            ]
            \addplot [draw=red,thick,name path=wool_constraint] {150-x/2}; \addlegendentry{Wool constraint: $3x + 6y \leq 900$}
            \addplot [draw=blue,thick,name path=cotton_constraint] {300-5*x/4}; \addlegendentry{Cotton constraint: $5x + 4y \leq 1200$}
            \addplot [draw=green, thick, name path=profit] {(4800-20*x)/16}; \addlegendentry{Profit: $20x + 16y = 4800$}
            \path[name path=xaxis] (0,0) -- (\xmax,0);
            \addplot[cyan!20] fill between[of=xaxis and wool_constraint,soft clip={domain=0:200}];
            \addplot[cyan!20] fill between[of=xaxis and cotton_constraint,soft clip={domain=200:\xmax}];
        \end{axis}
    \end{tikzpicture}
    \caption{Graphical Solution for Question 1c}
\end{figure}

\paragraph*{d)}
If 300 additional pounds of wool can be obtained, the wool constraint becomes $3x + 6y \leq 1200$. As can be seen in Figure 4, the optimal solution is on the point (133.33, 133.33).
On this point, the profit is $12(133.33) + 16(133.33) = 3733.3$.

\begin{figure}[H]
    \centering
    \begin{tikzpicture}
        \begin{axis}[
                legend style={nodes={scale=0.8}},
                grid,
                grid style={densely dashed},
                axis line style={->},
                xlabel=$x$,
                ylabel=$y$,
                xmin=0, xmax=\xmax,
                ymin=0, ymax=400,
                axis lines=center,
                axis on top=true,
                domain=0:\xmax,
            ]
            \addplot [draw=red,thick,name path=wool_constraint] {200-x/2}; \addlegendentry{Wool constraint: $3x + 6y \leq 1200$}
            \addplot [draw=blue,thick,name path=cotton_constraint] {300-5*x/4}; \addlegendentry{Cotton constraint: $5x + 4y \leq 1200$}
            \addplot [draw=green, thick, name path=profit] {(3733.3-12*x)/16}; \addlegendentry{Profit: $12x + 16y = 3733.3$}
            \path[name path=xaxis] (0,0) -- (\xmax,0);
            \addplot[cyan!20] fill between[of=xaxis and wool_constraint,soft clip={domain=0:133.33333333}];
            \addplot[cyan!20] fill between[of=xaxis and cotton_constraint,soft clip={domain=133.33333333:\xmax}];
        \end{axis}
    \end{tikzpicture}
    \caption{Graphical Solution for Question 1d}
\end{figure}

\section*{Question 2}
\paragraph*{a)}
The products should be produced as follows in order to minimize the total cost:
\begin{figure}[H]
    \includegraphics{Picture1.png}
\end{figure}
The minimized total cost is 32000\$.

\paragraph*{b)}
While the cost of producing cheesecake in machine 2 is in the range\\$[0, 4]$, the current basis remains optimal.

If the cost of producing cheesecake in machine 2 is increased by 1 unit, current solution does not change. Minimized cost becomes 35000\$.

\paragraph*{c)}
While the cost of producing cake in machine 2 is in the range $[4, \infty)$, the current basis remains optimal.

If the cost of producing cake in machine 2 is increased by 1 unit, current basis still remains optimal. In this case, minimized total cost is 32000\$.

If the cost of producing cake in machine 2 is decreased by two, current basis no longer remains optimal. In this case, we need to solve the problem again.
While solved again, the products should be produced as follows in order to minimize the total cost:
\begin{figure}[H]
    \includegraphics{Picture6.png}
\end{figure}
The minimized total cost is 30500\$.

\paragraph*{d)}
While the production requirement for muffins is in the range $[0, 7500]$, the current basis remains optimal.

The shadow price for this constrain is 3\$.
If the production requirement for muffins is increased by 1 unit, total cost increases by 3\$.

\paragraph*{e)}
While the time availability for machine 3 is in the range $[1000, \infty)$, the current basis remains optimal.

The shadow price for this constraing is 0\$.
If the time availability for machine 3 is changed, total cost does not change.

\paragraph*{f}
If the model is forced to produce at least one handmade cheesecake, the problem should be solved again.
While solved again, the products should be produced as follows in order to minimize the total cost:
\begin{figure}[H]
    \includegraphics{Picture7.png}
\end{figure}
However since this is a real life problem, counts must be integers. Therefore, if we round the non-integer values to the nearest integer, making sure that none of the constraints are violated, the products should be produced as follows in order to minimize the total cost:
\begin{figure}[H]
    \includegraphics{Picture9.png}
\end{figure}
In this case, the minimized total cost is 32002\$.

While the cost of producting a handmade cheesecake is in the range $[3, \infty)$, the current basis remains optimal.

For all values of the cost of producing a handmade cheesecake, handmade cheesecake takes a positive value because in this problem it is always greater than 1.
\end{document}
